\section{PDBx/mmCIF CA-contact target worked row}

This example uses the committed GAS target fixture. It is a target-normalization row, not an AOD prediction row.

\subsection{From atom-site rows to a contact bit}

The target derivation path is
\[
\texttt{mmCIF atom-site rows}
\to
\texttt{residue coordinate table}
\to
\texttt{CA distance}
\to
\texttt{contact bit}.
\]
For the pair $(1,3)$, the derived CA distance is
\[
 d_{1,3}=7.6\,\text{Angstrom}.
\]
The declared target rule uses threshold $8.0\,\text{Angstrom}$ and minimum sequence separation $2$:
\[
 d_{1,3}\le 8.0\,\text{Angstrom},
\qquad
 |3-1|\ge 2.
\]
Therefore the target contact bit is
\[
\boxed{O_{1,3}=1.}
\]
Adjacent pairs are retained in the distance ledger but are outside this contact-map scope.

\subsection{Ledger pointer}

The distance row is recorded in
\begin{center}
\small\path{manual-2/data/protein/pdb_mmcif_distance_matrix_derived.csv}
\end{center}
The contact row is recorded in
\begin{center}
\small\path{manual-2/data/protein/pdb_mmcif_contact_map_derived.csv}
\end{center}
The derivation manifest is
\begin{center}
\small\path{manual-2/data/protein/pdb_mmcif_contact_derivation_manifest.json}
\end{center}

\aodnoteblock{Target lock}{The target contact bit is forbidden as a raw D.E.C. premise and forbidden as a prediction-freeze premise. It is read only after the AOD packet is frozen.}
